% !TeX program = xelatex
% 《感知与人机交互》课程大作业报告 LaTeX 模板
% 编译方式：XeLaTeX

\documentclass[UTF8,a4paper,12pt,fontset=none]{ctexart}

\usepackage[left=1.30cm,right=1.30cm,top=2.95cm,bottom=2.00cm]{geometry}
\usepackage{fontspec}
\usepackage{xeCJK}
\usepackage{tikz}
\usepackage{eso-pic}
\usepackage{booktabs}
\usepackage{array}
\usepackage{caption}
\usepackage{titlesec}
\usepackage{setspace}
\usepackage{ragged2e}

% ==================== 字体 ====================
% AutoFakeBold 用于让宋体、楷体支持加粗
\setmainfont{Times New Roman}

\setCJKmainfont[
  AutoFakeBold=2.5
]{AR PL SungtiL GB}

\newCJKfontfamily\Song[
  AutoFakeBold=2.5
]{AR PL SungtiL GB}

\newCJKfontfamily\Kai[
  AutoFakeBold=2.5
]{AR PL KaitiM GB}

% ==================== 可编辑信息 ====================
\newcommand{\ReportName}{实验五：多层柜子中的物品类别识别}
\newcommand{\ReviewYear}{2026}

\newcommand{\StudentNameA}{王昊然}
\newcommand{\StudentIdA}{08123217}
\newcommand{\StudentScoreA}{}

\newcommand{\StudentNameB}{唐一宁}
\newcommand{\StudentIdB}{}
\newcommand{\StudentScoreB}{}

\newcommand{\StudentNameC}{方舟}
\newcommand{\StudentIdC}{}
\newcommand{\StudentScoreC}{}


% ==================== 正文格式 ====================
\setlength{\parindent}{2em}
\setlength{\parskip}{0pt}
\onehalfspacing

% 一级标题：楷体、加粗
\titleformat{\section}
  {\Kai\bfseries\zihao{-3}}
  {\thesection.}
  {0.45em}
  {}

\titlespacing*{\section}
  {0pt}
  {1.35em}
  {0.75em}

% 二级标题：楷体、加粗
\titleformat{\subsection}
  {\Kai\bfseries\zihao{4}}
  {\thesubsection}
  {0.55em}
  {}

\titlespacing*{\subsection}
  {2em}
  {1.00em}
  {0.45em}

\captionsetup[table]{
  name=表,
  labelsep=space,
  justification=centering,
  font={small},
  skip=5pt
}

\pagestyle{empty}

% ==================== 全文页眉和页码 ====================
\AddToShipoutPictureFG{%
  \begin{tikzpicture}[remember picture,overlay]

    % 封面页页眉加粗，正文页保持正常
    \ifnum\value{page}=0
      \node[
        anchor=center,
        font=\Song\bfseries\zihao{-5}
      ]
      at ([xshift=105mm,yshift=-18.2mm]current page.north west)
      {课程大作业};
    \else
      \node[
        anchor=center,
        font=\Song\zihao{-5}
      ]
      at ([xshift=105mm,yshift=-18.2mm]current page.north west)
      {课程大作业};
    \fi

    \draw[line width=0.60pt]
      ([xshift=20mm,yshift=-20.9mm]current page.north west)
      --
      ([xshift=190mm,yshift=-20.9mm]current page.north west);

    % 封面页页码加粗，正文页保持正常
    \ifnum\value{page}=0
      \node[
        anchor=center,
        font=\Song\bfseries\zihao{-5}
      ]
      at ([xshift=105mm,yshift=-277.0mm]current page.north west)
      {\thepage};
    \else
      \node[
        anchor=center,
        font=\Song\zihao{-5}
      ]
      at ([xshift=105mm,yshift=-277.0mm]current page.north west)
      {\thepage};
    \fi

  \end{tikzpicture}%
}

% ==================== 封面学生信息 ====================
% 整行以页面中心为基准
% 姓名、学号、评定成绩均加粗
\newcommand{\StudentLine}[3]{%
  \makebox[168mm][c]{%
    {\Song\bfseries 姓名：}%
    \underline{\makebox[34mm][c]{\Song\bfseries #1}}%
    \hspace{5mm}%
    {\Song\bfseries 学号：}%
    \underline{\makebox[38mm][c]{\Song\bfseries #2}}%
    \hspace{5mm}%
    {\Song\bfseries 评定成绩：}%
    \underline{\makebox[34mm][c]{\Song\bfseries #3}}%
  }%
}

\begin{document}

\Song
\setcounter{page}{0}

% ==================================================
% 封面页
% 所有文字均加粗
% ==================================================
\null

\begin{tikzpicture}[remember picture,overlay]

  \node[
    anchor=center,
    font=\Song\bfseries\fontsize{24pt}{30pt}\selectfont
  ]
  at ([xshift=105mm,yshift=-51mm]current page.north west)
  {东南大学};

  \node[
    anchor=center,
    font=\Kai\bfseries\fontsize{26pt}{34pt}\selectfont
  ]
  at ([xshift=105mm,yshift=-86mm]current page.north west)
  {《感知与人机交互》};

  \node[
    anchor=center,
    font=\Kai\bfseries\fontsize{26pt}{34pt}\selectfont
  ]
  at ([xshift=105mm,yshift=-114mm]current page.north west)
  {课程大作业};

  \node[
    anchor=center,
    font=\Kai\bfseries\fontsize{16pt}{21pt}\selectfont
  ]
  at ([xshift=105mm,yshift=-139mm]current page.north west)
  {名称：\ReportName};

  \node[
    anchor=center,
    font=\Song\bfseries\fontsize{14pt}{18pt}\selectfont
  ]
  at ([xshift=105mm,yshift=-164mm]current page.north west)
  {\StudentLine{\StudentNameA}{\StudentIdA}{\StudentScoreA}};

  \node[
    anchor=center,
    font=\Song\bfseries\fontsize{14pt}{18pt}\selectfont
  ]
  at ([xshift=105mm,yshift=-177mm]current page.north west)
  {\StudentLine{\StudentNameB}{\StudentIdB}{\StudentScoreB}};

  \node[
    anchor=center,
    font=\Song\bfseries\fontsize{14pt}{18pt}\selectfont
  ]
  at ([xshift=105mm,yshift=-190mm]current page.north west)
  {\StudentLine{\StudentNameC}{\StudentIdC}{\StudentScoreC}};

  \node[
    anchor=center,
    font=\Song\bfseries\fontsize{14pt}{18pt}\selectfont
  ]
  at ([xshift=105mm,yshift=-232mm]current page.north west)
  {%
    \makebox[168mm][l]{%
      评语：%
      \underline{\makebox[138mm][l]{}}%
    }%
  };

  \node[
    anchor=center,
    font=\Song\bfseries\fontsize{14pt}{18pt}\selectfont
  ]
  at ([xshift=105mm,yshift=-248mm]current page.north west)
  {%
    审阅教师：%
    \underline{\makebox[42mm][c]{}}%
  };

  \node[
    anchor=center,
    font=\Kai\bfseries\fontsize{14pt}{18pt}\selectfont
  ]
  at ([xshift=105mm,yshift=-262mm]current page.north west)
  {\ReviewYear 年 6 月 18 日};

\end{tikzpicture}

\newpage

% ==================================================
% 正文页
% ==================================================

\section{总体完成情况（主笔人：）}

\begin{table}[htbp]
  \centering
  \caption{小组人员及分工}

  \renewcommand{\arraystretch}{1.32}

  \begin{tabular}{
    >{\raggedright\arraybackslash}p{0.15\textwidth}
    >{\raggedright\arraybackslash}p{0.14\textwidth}
    >{\raggedright\arraybackslash}p{0.63\textwidth}
  }
    \toprule

    姓名
    &
    学号
    &
    实际完成的设计/开发/实验工作，承担任务占比例（\%）
    \\

    \midrule

    \rule{0pt}{1.65em} 王昊然& 08123217 & \\
    \rule{0pt}{1.65em} 唐一宁& 081231 & \\
    \rule{0pt}{1.65em} 方舟& 081231 & \\

    \bottomrule
  \end{tabular}
\end{table}

\section{总体原理与设计}

\vspace{1.65cm}

\section{具体分工开发实现}

\subsection{}

\vspace{0.85cm}

\subsection{}

\vspace{0.85cm}

\section{实验评估}

\vspace{1.10cm}

\section{分析总结}

\newpage

\section{附录}

\end{document}

